\section{Recurrence at the critical scale}\label{sec:recurrence}
The goal of this section is to place every favorite in a small disk at infinitely
many times. We first localize one favorite at an independent geometric time,
then make that favorite unique by an injective four-step insertion. Uniqueness
allows the resulting fresh segment to be combined with an arbitrary past.

\subsection{A resolvent estimate and reversal at a favorite}

\begin{lemma}[A stopped resolvent bound]\label{lem:stopped-resolvent}
Let \(0<z<1\), \(N=(1-z)^{-1}\), and let \(\mu\) be a finite positive
measure on \(\mathbb Z^2\) of mass \(\beta\).  Suppose that
\begin{equation}
(I-zP)\mu=\alpha\delta_0-\nu,\qquad \nu\ge0,
       \qquad \alpha-\nu(\mathbb Z^2)=(1-z)\beta.                 \tag{1}
\end{equation}
There are universal constants \(c,C>0\) such that, whenever
\(2\le f\le N^{1/8}\),
\begin{equation}
\mu\left(B\left(0,C\frac{\sqrt N}{f}\right)\right)
       \ge c\beta\,\frac{1+\log f}{f^2}.                         \tag{2}
\end{equation}
\end{lemma}

\begin{proof}
Let \(Q_R=[-R,R]^2\cap\mathbb Z^2\) and set
\[
 w_R:=|Q_R|^{-1}{\bf1}_{Q_R}*{\bf1}_{-Q_R}.
\]
Then \(0\le w_R\le1\), \(w_R(0)=1\), and \(w_R\) is supported in
\(B(0,3R)\).  Put
\[
 G_z:=\sum_{n\ge0}z^nP^n,\qquad \phi_R:=G_z*w_R.
\]
Fourier inversion gives
\[
 \widehat{\phi_R}(\theta)=
 \frac{|\widehat{{\bf1}_{Q_R}}(\theta)|^2/|Q_R|}
      {1-z\widehat P(\theta)}\ge0.
\]
Consequently \(\phi_R(y)\le\phi_R(0)\) for every \(y\).  Pairing (1) with
\(\phi_R\) and using \((I-zP)\phi_R=w_R\) yields
\begin{equation}
\begin{split}
 \mu(w_R)
 &=\alpha\phi_R(0)-\nu(\phi_R)\\
 &\ge(\alpha-\nu(\mathbb Z^2))\phi_R(0)
   =(1-z)\beta\phi_R(0).                                      \tag{3}
\end{split}
\end{equation}
The elementary Fourier integral over
\(\{\lVert\theta\rVert\le c/R\}\), together with
\(1-z\widehat P(\theta)\asymp (1-z)+\lVert\theta\rVert^2\),
gives, for \(R^2\le N/4\),
\begin{equation}
\phi_R(0)\ge cR^2\left(1+\log\frac N{R^2}\right).              \tag{4}
\end{equation}
Since \(\mu(w_R)\le\mu(B(0,3R))\), take
\(R=\lfloor\sqrt N/(4f)\rfloor\) in (3)--(4).  This proves (2),
after changing the absolute radius constant.
\end{proof}

\begin{lemma}[Localization of a selected favorite]\label{lem:plateau-bijection}
Let \(K\) be independent of the walk with
\[
       P(K=n)=(1-z)z^n,\qquad n\ge0.
\]
Attach independent uniform $[0,1]$ priorities \((U_x)_{x\in\mathbb Z^2}\) to
the sites and let \(X_K^*\) be the least-priority member of
\(\mathcal F_K\). Then
\begin{equation}
P\left(\lVert X_K^*\rVert_2\le C\frac{\sqrt K}{f}\right)
       \ge c\,\frac{1+\log f}{f^2},\qquad
       2\le f\le N^{1/8}.                                     \tag{5}
\end{equation}
More precisely, the proof gives a stopped-resolvent representation of the
endpoint law of \(-X_K^*\). The representation remains a positive mixture
satisfying (1) after restriction by any event measurable with respect to
the base tuple defined there.
\end{lemma}

\begin{proof}
We give the atomwise bijection; this is needed because a mere informal path
reversal does not control plateaux of tied favorites.

Fix a clock atom \(K=n\), a path
\(s=(s_0,\ldots,s_n)\), and first insert
\[
       1=\sum_x{\bf1}_{\{x=X_n^*\}}.
\]
Thus, in the sector currently considered, \(x\) is deterministic while the
priorities are translated.  Put \(m=M_n\), \(r=m-1\), and let
\(\sigma,\ell\) be the first and last visits to \(x\).  Reverse and center
the prefix through \(\ell\):
\[
       A_i=s_{\ell-i}-x,\qquad 0\le i\le\ell.
\]
It has exactly \(r\) positive returns to zero by time
\(\ell-\sigma\).  Hence it decomposes uniquely as
\begin{equation}
A=E_1\oplus\cdots\oplus E_r\oplus C,                    \tag{6}
\end{equation}
where the \(E_i\)'s are completed positive root excursions and \(C\)
starts at zero, has no positive return to zero, and ends at
\(a=-x\).  The centered forward suffix
\[
       B_i=s_{\ell+i}-x,\qquad 0\le i\le n-\ell,
\]
also has no positive return to zero.

Write \(H_r(y)\) for the number of visits to \(y\ne0\) in the completed
excursions and \(N_C(y),N_B(y)\) for the corresponding positive-time
counts in \(C,B\).  Translate priorities by the fixed sector shift,
\(V_y=U_{x+y}\).  The statement that the root is the priority-selected
favorite is exactly
\begin{equation}
H_r(y)+N_C(y)+N_B(y)
       \le r+{\bf1}_{\{V_y>V_0\}},\qquad y\ne0.                 \tag{7}
\end{equation}
Indeed, a better-priority site must be strictly below the root's value
\(r+1\), while a worse-priority site may tie it.

Expose the base tuple
\[
       {\cal B}=(r,E_1,\ldots,E_r,B,V)
\]
and define its fixed capacities
\begin{equation}
c_{\cal B}(y):=
       r+{\bf1}_{\{V_y>V_0\}}-H_r(y)-N_B(y),\qquad y\ne0.       \tag{8}
\end{equation}
Let \(\tau_c\) be the first positive return of a new walk \(C\) to zero or
the first time some \(N_C(y)\) exceeds (8).  The possible terminal arms
are exactly the prefixes \(C_0,\ldots,C_q\) with \(q<\tau_c\).

Conversely, a feasible base tuple and one such prefix reconstruct the
original path: concatenate (6), reverse it, translate by
\(-a=x\), and append the translated suffix \(B\).  Inequality (7) makes
\(x\) a favorite and excludes every better-priority tie; avoidance of zero
by \(B\) makes \(\ell\) its last visit.  The endpoint \(a\) recovers the
sector \(x=-a\), so both the path and the priority translation are unique.
The forward and inverse maps are therefore genuine bijections, including
\(r=0\).

If \(\rho_r\), \(q\), and \(k\) are the lengths of the completed part,
\(C\), and \(B\), reversal and symmetry preserve the path weight and
\begin{equation}
z^{\rho_r}4^{-\rho_r}\,
 z^q4^{-q}\,(1-z)z^k4^{-k}
       =(1-z)z^n4^{-n}.                                      \tag{9}
\end{equation}
Thus no multiplicity or normalization factor is hidden.

For a fixed capacity let
\[
       g_c(u)=E_0\sum_{q<\tau_c}z^q{\bf1}_{\{C_q=u\}}.
\]
Direct telescoping gives
\begin{equation}
(I-zP)g_c=\delta_0-E_0[z^{\tau_c}\delta_{C_{\tau_c}}].  \tag{10}
\end{equation}
Mixing (10) with the positive base weights in (9), and optionally
restricting those base weights by a base-measurable event, proves (1).

It remains to make the radius self-normalized.  Whole-path reversal,
translation to the new starting point, and transport of the iid priorities
preserve the joint law and send
\[
       (\sigma,\ell)\longmapsto(n-\ell,n-\sigma).
\]
As \(\sigma\le\ell\),
\begin{equation}
E(n-\sigma)=E\ell\ge n/2,
       \qquad P(n-\sigma\ge n/4)\ge1/3.                        \tag{11}
\end{equation}
Also \(P(K\ge N/2)\ge1/4\) for all large \(N\).  Hence the restriction
\[
       \Gamma=\{n-\sigma\ge N/8\}
\]
has mixture mass \(\beta\ge1/12\).  In the tuple,
\[
       n-\sigma=(\ell-\sigma)+(n-\ell)=\rho_r+|B|,
\]
so \(\Gamma\) is base-measurable, as required above.  Apply Lemma~\ref{lem:stopped-resolvent} to this restricted mixture. It gives a disk of radius $C_1\sqrt N/f$; on $\Gamma$, this is at most $\sqrt8 C_1\sqrt K/f$. Absorb the fixed factor into the radius constant.  This proves (5).
\end{proof}

\subsection{Removing ties and absorbing the past}

\begin{lemma}[Localization with a unique favorite]\label{lem:singleton-geometric-localization}
For the same geometric clock, uniformly for \(2\le f\le N^{1/8}\),
\begin{equation}
P\left(\mathcal F_K=\{X\},\
        \lVert X\rVert_2\le C\frac{\sqrt K}{f}\right)
       \ge c\,\frac{1+\log f}{f^2}.                            \tag{12}
\end{equation}
The conclusion remains valid, up to a universal change of \(c\), after
restricting the output clock to
\begin{equation}
aN\le K\le JN,\qquad J=\log N,\qquad
       a=\frac c8\,\frac{1+\log f}{f^2}.                       \tag{13}
\end{equation}
\end{lemma}

\begin{proof}
For every atom counted in (5), insert immediately after the last visit to
its priority-selected favorite \(x\) the fixed loop
\begin{equation}
x,\ x+e_1,\ x,\ x+e_2,\ x.                             \tag{14}
\end{equation}
The local time at \(x\) rises by two, those at the two neighbors rise by
one, and all other local times are unchanged.  Thus \(x\) is the unique
favorite of the new path.  The old suffix avoids \(x\), so in the image the
last visit to the unique favorite is the final vertex of (14).  Removing
the four preceding prescribed steps is therefore the unique inverse.  On
geometrically weighted path atoms the image-to-source weight ratio is
\(z^4/4^4\), bounded below for \(z\) close to one.  This proves (12).

Now \(P(K<aN)\le2a\), while
\(P(K>JN)\le\exp(-J/2)\).  With the constant in (13) chosen after the
constant in (12), these discarded masses total at most half the right-hand
side of (12).  Four added steps are absorbed by changing the universal
constants.
\end{proof}

\begin{lemma}[Adding an inherited local-time profile]\label{lem:inherited-profile-dichotomy}
Let \(H\ge0\) be any finitely supported integer profile and let
\(\ell\ge0\) be a second integer profile with a unique maximizer \(y\).
Every maximizer of \(H+\ell\) lies in
\(\operatorname{supp}H\cup\{y\}\).
\end{lemma}

\begin{proof}
If \(y\) is not a maximizer of \(H+\ell\), then
\(\max(H+\ell)>(H+\ell)(y)\ge\ell(y)\).  Outside
\(\operatorname{supp}H\), the profile \(H+\ell\) equals \(\ell\) and is at
most \(\ell(y)\), so every maximizer lies in \(\operatorname{supp}H\).
If \(y\) is a maximizer, uniqueness of the maximum of \(\ell\) shows that
the only possible maximizer outside \(\operatorname{supp}H\) is \(y\).
\end{proof}

\subsection{Infinitely many successful blocks}

\begin{proposition}[Recurrence at the critical scale]\label{prop:critical-recurrence}
Almost surely,
\begin{equation}
\liminf_{n\to\infty}
       \frac{R_n^+}{\sqrt{n/\log n}}=0.                        \tag{15}
\end{equation}
Consequently, for every \(\gamma\le\tfrac12\),
\begin{equation}
\liminf_{n\to\infty}\frac{R_n^+}{a_\gamma(n)}=0.
                                                                    \tag{16}
\end{equation}
\end{proposition}

\begin{proof}
Fix \(\varepsilon>0\).  Choose a deterministic sequence
\[
       J_{k+1}=J_k+10\log J_k,\qquad s_k=\lfloor e^{J_k}\rfloor,
\]
and, for the fresh block beginning at \(s_k\), put
\[
       N_k=\frac{e^{J_{k+1}}}{J_{k+1}^2},\qquad
       f_k=C\varepsilon^{-1}\sqrt{J_{k+1}}.
\]
Apply Lemma~\ref{lem:singleton-geometric-localization} to the fresh increment
walk, with an independent geometric clock with parameter $1/N_k$.  It gives a
unique fresh favorite \(Y_k\) satisfying
\begin{equation}
\lVert Y_k\rVert\le\frac{\varepsilon}{8}
             \sqrt{\frac{K_k}{J_{k+1}}}                       \tag{17}
\end{equation}
with conditional probability at least
\begin{equation}
d_k=c_\varepsilon\frac{\log J_{k+1}}{J_{k+1}},          \tag{18}
\end{equation}
even after
\begin{equation}
\frac{d_k}{8}N_k\le K_k\le J_{k+1}N_k.                 \tag{19}
\end{equation}
Since $e^{J_{k+1}}=e^{J_k}J_k^{10}$ and $J_{k+1}\sim J_k$,
(18)--(19) give, uniformly over the clock restriction,
\[
 c'_{\varepsilon}e^{J_k}J_k^7\log J_k
 \le K_k\le\frac{e^{J_{k+1}}}{J_{k+1}}
\]
for all sufficiently large $k$.  Thus $K_k/s_k\to\infty$, while
$s_k+K_k<s_{k+1}$, so each witness time belongs to its fresh block.
The singleton estimate already includes the four inserted steps.

Let
\[
 H_k(x)=\sum_{j=0}^{s_k-1}{\bf1}_{\{S_j=x\}},\qquad
 \ell_k(z)=\sum_{i=0}^{K_k}
       {\bf1}_{\{S_{s_k+i}-S_{s_k}=z\}}.
\]
The old profile stops before $s_k$, and the fresh profile includes its
time-zero visit.  Hence, for every $x\in\mathbb Z^2$,
\[
 H_k(x)+\ell_k(x-S_{s_k})
 =\sum_{j=0}^{s_k-1}{\bf1}_{\{S_j=x\}}
   +\sum_{i=0}^{K_k}{\bf1}_{\{S_{s_k+i}=x\}}
 =L_{s_k+K_k}(x).
\]
By
Lemma~\ref{lem:inherited-profile-dichotomy}, at time
\(t_k=s_k+K_k\) every favorite of the full walk belongs either to the old
range or to \(S_{s_k}+Y_k\).

Part (2) of Lemma~\ref{lem:elementary-walk-estimates} and Borel--Cantelli give
\[
       \max_{j\le s_k}\lVert S_j\rVert
       \le \sqrt{s_k}\,J_k
\]
eventually.  The clock bounds above give
\[
 \frac{s_kJ_k^2J_{k+1}}{K_k}
 =O_{\varepsilon}\left(\frac{J_k^{-4}}{\log J_k}\right),
 \qquad
 \log(s_k+K_k)=J_{k+1}+O_{\varepsilon}(\log J_{k+1})
\]
uniformly over (19).  Consequently,
\begin{equation}
\sqrt{s_k}\,J_k=o\left(\sqrt{\frac{K_k}{J_{k+1}}}\right),
 \qquad \log(s_k+K_k)\sim J_{k+1},\qquad s_k=o(K_k).            \tag{20}
\end{equation}
Combining (17)--(20), every full favorite at \(t_k\) is in
\begin{equation}
B\left(0,\frac{\varepsilon}{2}
                    \sqrt{\frac{t_k}{\log t_k}}\right)        \tag{21}
\end{equation}
for all sufficiently large successful \(k\).

Moreover \(J_k\sim10k\log k\), and hence
\begin{equation}
\sum_kd_k=\infty.                                      \tag{22}
\end{equation}
The fresh increments occupy disjoint deterministic blocks.  On the product
space enlarged by the independent clocks and priorities, let \(E_k\)
be the fresh singleton-success event (17), with clock restriction (19).
Its conditional probability given the past at \(s_k\) is at least
\(d_k\).  Lévy's conditional Borel--Cantelli lemma and (22) imply that
\(E_k\) occurs infinitely often almost surely.  Independently, the
old-range bound and the deterministic comparisons (20) hold eventually
almost surely; on their intersection, every sufficiently late \(E_k\)
implies (21).  Thus (21) occurs infinitely often.  Fubini removes the
auxiliary clocks and priorities: for almost every walk path, their
probability of producing infinitely many witness times is one, so that
path itself has infinitely many such integer witness times.

Intersecting over rational \(\varepsilon>0\) proves (15).  Finally,
\[
       \frac{a_{1/2}(n)}{a_\gamma(n)}
       =(\log n)^{\gamma-1/2}\le1
\]
eventually for \(\gamma\le1/2\), which proves (16).
\end{proof}
